\section{Findings}

\subsection{Troubleshooting begins by producing the missing system}

Initial questions frequently describe a visible symptom without the network of versions, interfaces, resources, constraints, and evidence required to interpret it. Respondents reconstruct that network through successive requests. They ask which software generation is installed, whether an interface is licensed, which method and dependencies were deployed, how a plate or tip is represented, what liquid and volume are involved, which physical actions completed, and what would count as a valid result. The work does not merely retrieve missing facts. It establishes which facts belong to the problem.

This pattern is clearest in scheduling and AI-method discussions. The scheduling scenario initially specified devices, tasks, durations, transfer time, and arrival times. Replies introduced stability windows, interruptibility, permitted storage, terminal disposition, priorities, and failure policy. Seven of eight incomplete requirement classes were surfaced, although none received a scenario-specific value. One participant summarized the dependency directly: ``Information about constraints becomes critical.'' \cite{forum_scheduling_toy} The discussion improved the representation of the problem without producing a completed schedule benchmark.

The AI-method discussion followed a similar path. An opening formulation described a finite operation vocabulary, volume, position, speed, and hardware boundaries. Replies added physical and biological behavior, sample value, human observation, implementation context, third-party devices, data handling, collaboration, regulation, and wet-process optimization. A practitioner stated: ``The human operator still needs to feed back the AI.'' \cite{forum_ai_method} The exchange reframed method generation from command composition into an activity embedded in material, organizational, and evidentiary conditions.

Participants also negotiate terminology. A hardware-involved check may be called a unit test, functional test, integration test, or system test depending on the speaker's engineering tradition. The collaborative task is to identify the object, environment, acceptance rule, evidence, and claim behind the label. This articulation work creates a temporary common information space in which people from software, automation, vendor, and laboratory backgrounds can reason about the same incident.

\subsection{Repair reconciles software histories with irreversible physical action}

Partial execution makes the limits of software-centered recovery visible. In one incident, tip pickup and aspiration completed before a destination-state error prevented dispensing. The initiating concern was material rather than exception handling alone: ``I might have valuable liquid in my tips.'' \cite{forum_error_catching} The exception identified a software conflict, although recovery depended on the physical consequences: source volume had changed, and a restart could duplicate or invalidate actions. Participants discussed prevalidation and queued execution as prevention mechanisms, then distinguished failures knowable from shared state from backend or hardware failures visible only at runtime.

A second discussion considered return of material after an abort. A custom handler demonstrated return-to-source behavior for one bounded software-error scenario \cite{forum_return_volume}. Participants simultaneously emphasized that a hardware crash can make the same action unavailable or unsafe. Repair therefore required a scenario-specific account of which actions completed, which resources remained controllable, and whether the material could retain scientific validity.

Across the observability cases, final disposition was consistently weak. Logs often contained a warning, curve, or exception without a linked account of sample state, manual intervention, recovery action, or scientific outcome. Participants could identify a divergence and remain unable to reconstruct whether a sample was retained, repeated, or discarded. The strongest descriptive relationship in the coded register connected observability and recovery. The qualitative mechanism is relational: repair quality depends on how the community can connect command history, material identity, physical observation, modeled state, intervention, and result.

This repair work extends beyond the immediate run. A version-control incident revealed that a directory assumed to contain source included an actively maintained recovery workspace. Restoring files through source control did not restore an earlier system state because the vendor environment continued to update the workspace. Participants had to distinguish method source, deployment configuration, generated identity, physical teaching, and runtime recovery state. The repair involved revising the conceptual boundary around what version control should contain.

\subsection{Public support is porous and outcomes migrate}

The forum creates public visibility without containing the entire support process. Some discussions move toward vendor contact, direct messages, Slack groups, meetings, paid development kits, or local code review. Migration can accelerate resolution by bringing the right expertise or restricted information into the interaction. It also removes evidence from the public record and weakens later discoverability.

The labware-database discussion illustrates this tension. The initiating problem was stated simply: ``Finding labware is a mess.'' \cite{forum_labware_database} The discussion then expanded toward shared dimensions, tolerances, material, part numbers, and provenance. Participation grew, and coordination moved toward Slack and meetings. The migration may have supported collaboration, but the public thread no longer contained a complete account of subsequent decisions, schema changes, or maintenance ownership. The public community remained the entry point and lost the authoritative outcome.

Commercial support structures create another boundary. In a scheduler-driver discussion, detailed templates, examples, documentation, and developer support existed through a commercial development kit, while the source of existing drivers remained inaccessible. In another case, a driver could be installed but required an additional feature licence to operate. Participants discussed vendor routing and external serial wrappers; external control then raised a state-synchronization problem with the platform. The public discussion exposed the architecture and access conditions without fully replacing the restricted support channel.

Public outcome therefore has at least three components: whether the immediate problem was resolved, whether the evidence remained visible, and whether the solution became maintainable. A private answer can restore operation and fail as community infrastructure. An open thread can generate valuable requirements or diagnostic distinctions without a confirmed local resolution.

\subsection{Durable knowledge requires infrastructuring after diagnosis}

The strongest positive cases include work after the technical answer. A shared tip definition worked for its original contributor and failed for another user. The subsequent process identified several physical and coordinate causes, recorded dates and lessons, corrected the software definitions, merged changes, and published a post-mortem \cite{forum_tip_postmortem}. The incident became durable because participants connected failure evidence, correction, code, and explanation.

A public plate-reader release followed another durable pattern. It exposed a host-independent API, executable example, contribution path, and identifiable maintainer. The release did not claim complete feature coverage. Its value came from making the integration surface inspectable and extensible while retaining its limits.

Documentation discussions show how small interventions can produce infrastructure. In one case, reference material was already installed inside a vendor environment and remained undiscovered until a community answer identified it. The requester responded: ``This would have saved me so much time.'' \cite{forum_hsl_learning} In another, maintainers explained local documentation builds and a page-level edit pathway. These responses did more than answer an immediate question. They lowered the cost of future contribution and correction.

Durability also requires ownership and lifecycle. Participants repeatedly encountered advice that was version-specific, privately held, stale, or detached from its original evidence. The proposed resolved-knowledge record emerged from this pattern. It attaches applicability, root-cause status, validation stage, evidence grade, limitations, maintainer, last verification, source provenance, and supersession to an answer. The record is an infrastructuring proposal: it turns resolution into an object that can be maintained, disputed, and corrected.
